\documentclass{ctexbook}
\usepackage{amsmath}
\usepackage{bm}
\usepackage{amssymb}
\usepackage{boondox-cal}
\title{快乐的高中物理}
\begin{document}
	\chapter{电磁感应}
	\section{基础知识：法拉第电磁感应定律}
	\subsection{知识回顾：什么是磁通量}
	我们在之前的物理课中，给出了一个叫做磁通量的物理量$\varPhi$。这个物理量在接下来的课程中比较重要，我们再次强调一下。
	
	\paragraph{两特殊方向磁通量的计算}~{}
	\\
	\indent
	高中物理中，若磁感线与某个面垂直（与该面法向量平行），我们通过$\varPhi=BS$来计算，这个物理量用来描述穿过单位面积的磁感线的多少。
	
	两种极限情况：
	
	\hspace*{\fill}
	
	\hspace*{\fill}
	
	\hspace*{\fill}
	
	\hspace*{\fill}
	
	\hspace*{\fill}
		
		\,\,\,\,\,\,\,\,\,\,\,\,\,\,\,\,\,\,\,\,\,\,\,\,\,\,\,\,\,\,（$B\perp S$）\,\,\,\,\,\,\,\,\,\,\,\,\,\,\,\,\,\,\,\,\,\,\,\,\,\,\,\,\,\,\,\,\,\,\,\,\,\,\,\,\,\,\,\,\,\,\,\,\,\,\,\,\,\,\,\,\,\,\,\,\,\,\,\,\,\,\,\,\,\,（$B//S$）
	
	\paragraph{任意方向磁通量的计算}~{}
	\\
	\indent
	原理：磁感应强度是矢量$\vec{B}$，是可以被分解为平行于面的磁感强度$B_{//}$和垂直于面的磁感强度$B_{\perp}$。只有\underline{\hbox to 10mm{}}于面的磁感强度分量对磁通有贡献。即$\varPhi =$\underline{\hbox to 10mm{}}。
	
	Tips：不要记乘角度的某个三角函数，题目所给的角和你记忆的角度可能不一样。书上所说乘以$\cos \theta$可能是这道题的$\sin \alpha$。
	
	\hspace*{\fill}
	
	Another计算方法：磁通量更准确的定义是
	
	\centerline{$\varPhi=\vec{B} \cdot \vec{S}=\vec{B} \cdot \vec{n} S$}
	
	其中我们把面积矢量化，面积矢量的方向即为该面法向量方向，模长为该面面积大小。这样磁通量$\varPhi$即可看作磁感强度$\vec{B}$和面积矢量$\vec{S}$做点乘的结果。
	
	\paragraph{磁感强度具有数量可加性}~{}
	\\
	\indent
	磁感应强度是没方向的数量，且满足数量可加性。即若一个面$S$可分为两部分$S_1,S_2$，通过这两个面的磁通为$\varPhi_1,\varPhi_2$,那么$\varPhi=\varPhi_1+\varPhi_2$。同理，若分为$n$部分，则$\varPhi=\varPhi_1+\varPhi_2+\cdots+\varPhi_n$。
	
	思考题：圆环半径与磁场半径不同的磁通量计算
	
	\hspace*{\fill}
	
	\hspace*{\fill}
	
	\hspace*{\fill}
	
	\hspace*{\fill}
	
	\hspace*{\fill}
	
	\paragraph{易混！磁通和磁链的差别！}~{}
	\\
	\indent
	引入题：一匝半径为$r$，匝数为$N$圆形线圈，线圈内部有垂直于面的磁感线，磁感强度为$B$，求磁通量。
	
	答案：$\pi r^2 B$ or $N \pi r^2 B$
	
	正确答案是$\pi r^2 B$。How to 理解？平面只有一层！即若磁感线在穿过平面范围内是直的，磁感线只能穿过这个面一次，贡献一次磁通量。这样的话何来$N$？
	
	了解：事实上，后者是物理或工程中的另外一个物理量，叫做磁链$\varPsi$，且有这样的关系，$\varPsi=N\varPhi$。高中阶段没有给出这个物理量的定义，那么凡是匝数为$N$，选项中有$N$的直接排除！
	\subsection{法拉第电磁感应定律}
	
	在初中物理中，我们就学过电磁感应的现象。现在我们来定量描述这个物理现象。
	
	注意：高中物理是带领你描述各种各样的物理现象的，你只需要知道这个现象可以被这样描述，且能运用到题目中的模型中去即可！一定不要死纠为什么会产生这样的现象，为什么这个公式长成这个样子，你可以理解为前人通过实验或者理论推导得出的结果。在高中阶段若以提分为导向，不要去深究。若确实感兴趣，在保证高中做题不受影响的情况下，你可以通过互联网或者查阅大学专业教材获得更详细的解释。
	
	\paragraph{感应电动势大小计算}~{}
	\\
	\indent
	
	通过前人的实验研究测量总结，我们得到了关于定量描述法拉第电磁感应定律的式子：
	
	\centerline{$\mathcal{E}=\dfrac{\Delta \varPhi}{\Delta t}$}
	
	事实上，这只是计算平均感应电动势的式子，左侧的$\mathcal{E}$应记为$\overline{\mathcal{E}}$。若要计算瞬时值，上式要取$\Delta t\rightarrow 0$的极限，在学习导数后，我们发现感应电动势实际上是$\varPhi$对时间$t$求导。
	
	\centerline{$\mathcal{E}=\varPhi'_t=\dfrac{\mathrm{d}\varPhi}{\mathrm{d}t}$}
	
	\paragraph{动生电动势和感生电动势}~{}
	\\
	\indent
	
	$\mathcal{E}=\dfrac{\mathrm{d}\varPhi}{\mathrm{d}t}=\dfrac{\mathrm{d} (BS)}{\mathrm{d}t}=S\dfrac{\mathrm{d}B}{\mathrm{d}t}+B\dfrac{\mathrm{d}S}{\mathrm{d}t}=\mathcal{E}_	c+\mathcal{E}_M$
	
	我们看后一项，发现这一项产生的原因是\underline{\hbox to 30mm{}}，这通常意味着运动，故我们把这一项称为\underline{\hbox to 20mm{}}，且在后面的课程中，我们发现导致这个电动势产生的非静电力$K$是\underline{\hbox to 30mm{}}。
	
	我们看前一项，发现这一项产生的原因是\underline{\hbox to 30mm{}}，在后面的课程中，我们发现导致这个电动势产生的非静电力$K$是\underline{\hbox to 30mm{}}，故我们把这一项称为\underline{\hbox to 20mm{}}。
	
	\paragraph{导轨上金属棒$\mathcal{E}=Blv$的推导}~{}
	\\
	\indent
	
	\hspace*{\fill}
	
	\hspace*{\fill}
	
	\hspace*{\fill}
	
	\hspace*{\fill}
	
	\hspace*{\fill}
	
	\hspace*{\fill}
	
	\hspace*{\fill}
		
	\paragraph{不同类型的动生电动势：垂直切割or不垂直切割，平动or转动，有效距离的判断and计算}~{}
	\\
	\indent
	
	\hspace*{\fill}
	
	\hspace*{\fill}
	
	\hspace*{\fill}
	
	\hspace*{\fill}
	
	\hspace*{\fill}
	
	\hspace*{\fill}
	
	\hspace*{\fill}
	
	\hspace*{\fill}
	
	\hspace*{\fill}
	
	\hspace*{\fill}
	
	\hspace*{\fill}
	
	\hspace*{\fill}
	
	\hspace*{\fill}
	
	\hspace*{\fill}
	
	\hspace*{\fill}
	
	\hspace*{\fill}
	
	\hspace*{\fill}
	
	\hspace*{\fill}
	
	\hspace*{\fill}
	
	\hspace*{\fill}
	
	\hspace*{\fill}
	
	\hspace*{\fill}
	
	\hspace*{\fill}
	
	\hspace*{\fill}
	
	\hspace*{\fill}
	\paragraph{感应电动势方向判断：楞次定律与右手定则}~{}
	\\
	\indent
	
	在上面的讨论中，我们知道了如何计算感应电动势的大小，那我们如何来判断电流方向？
	
	楞次定律：感应电流所产生的一切后果，必将反抗引起它的所有原因。
	
	图解：
	\hspace*{\fill}
	
	\hspace*{\fill}
	
	\hspace*{\fill}
	
	\hspace*{\fill}
	
	\hspace*{\fill}
	
	右手定则：\underline{\hbox to 103mm{}}
	
	\underline{\hbox to 120mm{}}
	
	\subsection{一些物理学史}
	发现电磁感应现象（1831）：法拉第
	
	总结楞次定律（1834）：楞次
	
	总结出法拉第电磁感应定律（1845、1846）：纽曼、韦伯
	
	\newpage
	
	\section{题型1：电磁感应在电路中的应用}
	
	电磁感应的重大作用即为提供电动势，除开极少数考察概念的题目，所有的电磁感应题目都是在电路的大背景内考察。因此熟练掌握电磁感应在电路中的应用，是解决一切电磁感应题目的基础。
	
	在引入电磁感应之前，我们需要先复习一些电路中的知识点。
	
	\subsection{知识回顾：电路}
	
	基本电路：
	
	\hspace*{\fill}
	
	\hspace*{\fill}
	
	\hspace*{\fill}
	
	\hspace*{\fill}
	
	\hspace*{\fill}
	
	直流电路由\underline{\hbox to 20mm{}}和\underline{\hbox to 20mm{}}构成，电流恒定不变。
	
	电源是通过\underline{\hbox to 20mm{}}做功把其他形式能量转化为电能的装置。
	
	下面我们着重讨论电源的电学参量。我们常常用电动势和内阻来刻画一个电源。
	
	在电路中，我们有闭合电路欧姆定律：
	
	\centerline{$\mathcal{E}=Ir+IR$}
	
	除此之外，还有路端电压这个重要的概念。事实上，$U_ab=IR$常常不会以路端电压这个名次考察，而是以让你计算两点的电压（电势差）的形式考察。
	
	\subsection{电磁感应中的电路}
	由法拉第电磁定律，我们知道了变化的磁通量可以产生相应的电动势。不难想到，参与电磁感应的部分充当电路的电源部分。
	
	电磁感应中常考的作为电源的装置是金属杆，而在选择题中，充当电源的部分可能是三角形，弧线各种不规则形状的导线。
	
	分辨出哪部分充当电源，哪部分充当负载是电磁感应中最重要的步骤！之后我们联系上电磁感应和电路的知识列方程求解即可。
	
	\subsection{一个重要的结论：$q=n\dfrac{\Delta \varPhi}{R_{\mathbf{sum}}}$}
	
	这个结论不能直接在大题中使用，需要推导后使用。
	
	推导：$q=\overline{I}\Delta t=\dfrac{\overline{\mathcal{E}}}{R_{\mathbf{sum}}}\Delta t=\dfrac{n \dfrac{ \Delta \varPhi}{\Delta t}}{R_{\mathbf{sum}}}\Delta t=n\dfrac{\Delta \varPhi}{R_{\mathbf{sum}}}$
	
	注意：
	
	1.此处的$R_{\mathbf{sum}}$是指整个电路中的总电阻，既包括电源内阻，又包括电路上的负载。
	
	2.此处的电荷是指流过电路的净电荷，即$q=0$不代表电路中无电流通过。
	
	\newpage
	
	\section{题型2：电磁感应中的力学问题（R-杆电路）}
	
	根据我们前面的电路分析，我们应该要得到各元件的各种电学参量。这时由于电流的存在，在磁场中的物体将受到安培力的作用，从而改变运动状态。而运动状态的改变则会使感应电动势得到改变，进一步改变电学参量。即整个过程中力学量和电学量在相互改变，最后达到动态平衡。这一讲我们将会通过几组模型来研究这个过程。
	
	电磁感应中最简单的模型是不同的单杆模型，当然也有更难的双杆问题，我们将一一列举，并分析电和力是如何相辅相成融为一体的。
	
	在后面，我们会看到电磁感应与动力学的问题中磁场的方向对运动状态不会有影响。
	
	\subsection{牛顿力学分析}
	
	这是最简单的电路，即一个电阻和一个在导体中切割的金属杆组成。
	
	R-杆基本电路：
	
	\hspace*{\fill}
	
	\hspace*{\fill}
	
	\hspace*{\fill}
	
	\hspace*{\fill}
	
	\hspace*{\fill}
	
	这个电路又可以分为有无外力和有无初速度。统一起见，我们先设外力为$F$，初速度为$v_0$。无外力或无初速度的情况即为相应物理量为0的解。
	
	我们先进行电路分析，每一时刻的速度为$v$，即$v=v(t)$。则电动势$\mathcal{E}=BLv$。据闭合电路欧姆定律，有电流$I=\dfrac{\mathcal{E}}{R+r}=\dfrac{BLv}{R+r}$。此时，导体杆受到的安培力为$F_A=BIL=\dfrac{B^2L^2v}{R+r}$,方向由右手定则给出，可发现与运动方向相反，即向左。这个的理解是显然的，可以通过楞次定律去解释它。
	
	我们对这个杆受力分析列牛顿第二定律的动力学方程，则有：
	
	\centerline{$ma=m\dfrac{\mathrm{d}v}{\mathrm{d}t}=F-F_A=F-\dfrac{B^2L^2}{R+r} v$}
	
	这是一个一阶线性微分方程，在高中阶段不要求求解出具体的函数表达式。但我们不难通过观察发现，$v=\dfrac{F(R+r)}{B^2L^2}$是这个方程的平凡解，将此解代入方程后，发现$a=\dfrac{\mathrm{d}v}{\mathrm{d}t} \equiv 0$，即速度不会随着时间改变。即如果速度和外力满足上式，金属杆将会一直保持着这个速度匀速运动下去。我们把这个速度记为$v_C$。
	
	如果$v<v_C$，式子右边计算出正数项，即$a=\dfrac{\mathrm{d}v}{\mathrm{d}t} > 0$，速度将增大。但随着速度增大，式子右边的正数项数值上将减少，即加速度在数值上减少，速度增加得更慢。整体来看，金属杆将做加速度减少的加速运动。
	
	反之，若$v>v_C$，同上分析，金属杆将做加速度减少的减速运动。
	
	那金属杆初速度小于$v_C$时，能不能被外力加速到超过$v_C$呢？即是，在这样一个动力学系统中，其$v-t$图像中的任意曲线能否跨过$v=v_C$这条直线呢？
	
	答案是否定的，这是通过求解上面微分方程所得到的。这里我们不加推导的给出上述方程的非平凡解。
	
	\centerline{$v(t)=v_C+(v_0-v_C)e^{-\frac{B^2 L^2}{m(R+r)} t}$}
	
	当$t \rightarrow \infty $时，$v \rightarrow v_C$，但永远不会跨过$v=v_C$这条直线！！
	
	试做出下列条件的$v-t$图，且标记出当$t \rightarrow \infty $时的渐进行为，有多种可能结果的要分类讨论：
	
	（1）$v_0=\dfrac{I}{m}$，$F=0$；
	
	（2）$v_0=0$，$F=mg \sin \theta$；
	
	（3）$v_0=v_0$，$F=mg \sin \theta$。
	\hspace*{\fill}
	
	\hspace*{\fill}
	
	\hspace*{\fill}
	
	\hspace*{\fill}
	
	\hspace*{\fill}
	
	\hspace*{\fill}
	
	\hspace*{\fill}
	
	\hspace*{\fill}
	
	\hspace*{\fill}
	
	\hspace*{\fill}
	
	\hspace*{\fill}
	
	\hspace*{\fill}
	
	\hspace*{\fill}
	
	\hspace*{\fill}
	
	\hspace*{\fill}
	
	\hspace*{\fill}
	
	\hspace*{\fill}
	
	\hspace*{\fill}
	
	\hspace*{\fill}
	
	\hspace*{\fill}
	
	\hspace*{\fill}
	
	\hspace*{\fill}
	
	\hspace*{\fill}
	
	\hspace*{\fill}
	
	\hspace*{\fill}
	
	\hspace*{\fill}
	
	\hspace*{\fill}
	
	\hspace*{\fill}
	
	\hspace*{\fill}
	
	\hspace*{\fill}
	
	从画图中，我们可以发现，当$t$较大后，速度$v$将接近与$v_C$。由于现实中其他涨落的影响，速度从观测的角度而言可认为恒定不变，这时候我们说系统达到收尾速度$v_C$，我们将经常在题目中看到金属杆最后匀速运动，其实就是说的是$v$近似等于$v_C$这个过程。若题目中要求计算这个匀速运动的速度，那计算$v_C$即可。
	
	\subsection{能量分析}
	
	以上我们是从牛顿第二定律的角度分析得到的，下面我们从能量和做功的方面去考虑。
	
	\centerline{$W_F+W_{F_A}=\dfrac{1}{2}m v_C^2-\dfrac{1}{2}m v_0^2$}
	
	我们可以把外力$F$做功视为金属杆$F$势场中的势能转化为其他形式的能。这里其他形式的能量不仅包含速度改变的动能，也包括电流在电路中所产生的焦耳热。对应到上式，得到如下功能关系
	
	\centerline{$-W_{F_A}=E_{Circuit}=Q_{Joule}$}
	
	即，克服安培力做的功等于电路中产生的电能。若电路为纯电阻电路，上式第二个等号成立，即克服安培力做的功等于电路中产生的焦耳热。
	
	\subsection{动量分析}
	
	这个系统在动量这里没有什么特别的，不用怎么去分析。
	
	\section{题型2：电磁感应中的力学问题（E-杆电路）}
	在考试中，单纯的R-杆电路不可能用来考整个大题，因为这样一个指数衰减向特定速度的且伴随着能量耗散运动过程是无趣的，我们希望有更多有意思的运动，这需要我们加入一个外界电源来驱动金属杆运动。这样的电源可以是我们的直流电源，也可以是我们的另一个金属杆运动引起的电动势。这就引出了这节所讲的两种电路。
	
	下面这节我们先介绍第一种电路。
	
	\subsection{基本分析}
	
	学这种电路是我们所说学后面双杆电路的基础，我们要认真学懂这个电路。
	
	E-杆基本电路：
	
	\hspace*{\fill}
	
	\hspace*{\fill}
	
	\hspace*{\fill}
	
	\hspace*{\fill}
	
	\hspace*{\fill}
	
	注意，我们呈现的是基本的电路，题目的电路可能还会对这样的电路进行改造，但掌握本质我们便能解决这些题。
	
	电路接通后，由于存在电动势$\mathcal{E}$，电路中将会有电流$I=\dfrac{\mathcal{E}}{R}$。这样的话我们的金属棒将会受到安培力的作用，$F_A=BIL$。我们的金属棒由于受到安培力就会开始运动。而金属棒一旦开始运动，就会开始切割磁感线造成电磁感应。可以判断，电磁感应的电动势将会影响直流电源的电动势作用，这样的话电流就会减少，安培力也会减少，加速度将逐渐减慢。用方程描述这一过程则为
	
	\centerline{$ma=m\dfrac{\mathrm{d}v}{\mathrm{d}t}=BL\dfrac{\mathcal{E}-BLv}{R}$}
	
	这个方程和上一节课所列出的方程相似度如此之高，以至于我们可以大胆猜想它的解和上一节的解形式类似，都是e指数变化到所谓的收尾速度$v_C$。事实上，这样的猜想是正确的，并且我们可以显然的给出$v_C=\dfrac{\mathcal{E}}{BL}$，这不过就是一个看出方程平凡解的过程。
	
	这道题的另一种理解过程就是把上一节R-杆电路中的外力$F$用直流电动势所造成的安培力替代，即
	
	\centerline{$F=\dfrac{BL\mathcal{E}}{R}$}
	
	此时，我们可以完全套用上节课所有的结论。
	
	\subsection{类比分析：能量怎么转换的}
	
	由于我们已经分析出E-杆电路和R-杆电路在本质上都是同一个电路，我们不难直接利用上一节的结论。
	
	我们先列出此时外力做的功$W_F=\dfrac{BLx\mathcal{E}}{R}=\dfrac{B \Delta S}{R} \mathcal{E}=q \mathcal{E}$，这和我们直接用电动势的定义得到的结果是一样的。
	
	整个过程即为电动势的能量给了金属棒使其有了动能，且让电路中有了电流，使整个系统有了焦耳热。
	
	\subsection{动量分析}
	这个电路的动量依然没什么好分析的。
	
	\newpage
	
	\section{题型2：电磁感应中的力学问题（杆-杆电路/双杆电路）}
	之前我们学的R-杆和E-杆模型都是只有一个杆，因此只有一个对象的动力学参量，比较方便。而在这节课，我们将学习双杆模型，它具有两个研究对象，即有两对动力学参量，相较而言会更加麻烦。但由于两个物体的引入，导致可以列出动量守恒或类似于动量守恒的式子，这个式子会大大简化这类问题的解答过程。因此，我们需要重点掌握动量这个方法来解决双杆问题。
	
	在此之前，我们还是先通过受力分析列出最基本的动力学方程。
	
	\subsection{基本分析：明确运动趋势}
	
	杆-杆基本电路：
	
	\hspace*{\fill}
	
	\hspace*{\fill}
	
	\hspace*{\fill}
	
	\hspace*{\fill}
	
	\hspace*{\fill}
	
	\hspace*{\fill}
	
	\hspace*{\fill}
	
	\hspace*{\fill}
	
	显然，我们可以把其中的一个杆等效为一个可变电动势的电源，这样我们即把双杆问题退化为一个E-杆问题和一个确定电动势的问题。显然，E-杆问题对我们而言是一个已解决的问题，而确定电动势的问题也通过法拉第电磁感应定律解决。就这样，我们把一个我们不知道的问题拆开，转化为两个我们已经解决的问题。
	
	下面尝试用牛顿第二定律列出动力学方程。这里有两个杆，显然要列出两个方程。我们规定，主动杆为2杆，被动杆为1杆（即是我们规定1杆上的外力$F_1 \equiv 0$）。
	
	两个方程：
	
	\hspace*{\fill}
	
	\hspace*{\fill}
	
	\hspace*{\fill}
	
	\hspace*{\fill}
	
	不难发现，这是一个一阶常微分方程组，可以大胆猜想，其解的形式也会同前两个问题一样以指数衰减的形式达到一个收尾速度。事实上，当主动杆外力为零是，解得形式确实是这样。在此我们不再给出这个解的形式，因为其对我们高中物理学习没有任何的帮助。
	
	下面，我们感兴趣的其实是收尾速度的大小。这样的速度会使得两杆产生等量的电动势，整个电路回路内将不会有电流。即
	
	\centerline{$B l_1 v_{C1}=B l_2 v_{C2}$}
	
	或两杆产生电动势的差是一个定值，即
	
	\centerline{$B l_1 v_{C1}-B l_2 v_{C2}=\mathrm{const}$}
	
	但这样我们只找到了一个方程，显然无法解除两个收尾速度。部分问题下，剩下的方程可以通过受力分析给出。其他的问题下，剩下的方程需要通过动量分析得出。
	
	\subsection{动量分析：双杆问题的关键}
	列出动量定理的式子：
	
	\centerline{$m \Delta v =F \Delta t=BIL \Delta t=BLq=\dfrac{B^2 L^2 v \Delta t}{R}=\dfrac{B^2 L^2 \Delta x}{R}$}
	
	显然，对两个杆分别列动量定理，在式子右侧将会有电荷量$q$出现。双杆问题都在一个回路中，因此两个式子的$q$相等，是可以约去的。这样就得到了第二个式子。
	
	\subsection{主动杆外力不为零？还是动量！}
	外力不为零意味着多了一个非齐次项，可以分析，这时候的两杆将不会有一个确定的收尾速度，而是有一个确定的收尾加速度。解决这样的问题的重点依然在对动量的分析上，只是会多出一个$F \Delta t$项。
\end{document}